\documentclass[11pt,a4paper]{article} % 11pt zamiast 12pt

\usepackage[T1]{fontenc}
\usepackage[utf8]{inputenc}
\usepackage[polish]{babel}
\usepackage[a4paper,margin=2cm]{geometry} % mniejsze marginesy

\usepackage{amsmath,amssymb}
\usepackage{listings}
\usepackage{enumitem}
\usepackage{titlesec}

% kompaktowe sekcje
\titlespacing*{\section}{0pt}{8pt}{4pt}
\titlespacing*{\subsection}{0pt}{6pt}{3pt}

% kompaktowe listy
\setlist[itemize]{noitemsep, topsep=2pt}
\setlist[enumerate]{noitemsep, topsep=2pt}

% kod bardziej zbity
\lstset{
    basicstyle=\ttfamily\footnotesize,
    breaklines=true,
    columns=fullflexible,
    aboveskip=6pt,
    belowskip=6pt
}

\linespread{0.96}

\title{\textbf{Sprawozdanie 1}\\Wybrane Zagadnienia Algebry}
\author{Jakub Kowal\\279706}
\date{}

\begin{document}
\maketitle

Użyte parametry: $a=2,\ b=7,\ c=9,\ d=7,\ e=0,\ f=6$.

\section*{1. $\mathbb{Z}[i]$}

\textbf{Reprezentacja:} \texttt{Pair(a,b)} oznacza $a+bi$ (\texttt{first} – rzeczywista, \texttt{second} – urojona).

\textbf{Norma:} $N(a+bi)=a^2+b^2$.

\textbf{Dzielenie:} przybliżenie ilorazu + sprawdzanie floor/ceil; poprawne gdy $N(r)<N(b)$.

{\footnotesize
\begin{lstlisting}
function Dziel_wszystkie(a::Pair{<:Integer,<:Integer},b::Pair{<:Integer,<:Integer})
 denominator = float(Norma(b))
 real = (a.first*b.first + a.second*b.second) / denominator
 imag = (a.second*b.first - a.first*b.second) / denominator

 potential_reals = Set([floor(Int, real), ceil(Int, real)])
 potential_imags = Set([floor(Int, imag), ceil(Int, imag)])

 wyniki = Vector{Tuple{Pair{Int,Int},Pair{Int,Int}}}()
 for r in potential_reals, i in potential_imags
  q = Pair(r, i)
  rem = Odejmij(a, Mnoz(b, q))
  if Norma(rem) < Norma(b)
   push!(wyniki, (q, rem))
  end
 end
 return wyniki
end
\end{lstlisting}
}

\textbf{NWD/NWW:} Euklides, $\mathrm{NWW}(a,b)=\frac{ab}{\mathrm{NWD}(a,b)}$.

{\footnotesize
\begin{lstlisting}
function NWD(a::Pair{<:Integer,<:Integer},b::Pair{<:Integer,<:Integer})
 while !(b.first == 0 && b.second == 0)
  _,r=Dziel(a,b)
  a=b; b=r
 end
 return a
end
\end{lstlisting}
}

\textbf{Wyniki:}
$N(2-7i)=53$.

$\frac{11+14i}{6i}$ daje:
\begin{enumerate}
\item $2-i,\ 5+2i$
\item $2-2i,\ -1+2i$
\item $3-3i,\ -1-4i$
\end{enumerate}

NWD$=1$, NWW$=-217+539i$ (do mnożenia przez $\{1,-1,i,-i\}$).

{\footnotesize
\begin{lstlisting}
NWD([a+bi]) = 2 + 7i
NWD([])     = 0+0i
\end{lstlisting}
}

\section*{2. $\mathbb{R}[x]$}

\textbf{Reprezentacja:}
{\footnotesize
\begin{lstlisting}
struct Pierscien{T}
 coeffs::Dict{Vector{Int}, T}
end
\end{lstlisting}
}

\textbf{Operacje:} norma = stopień, dzielenie klasyczne, NWD – Euklides, NWW – iloczyn/NWD.

\textbf{Wyniki:}

Norma: $2$.

\[
\frac{9x^2+7}{x+1}=9x-9,\quad r=16
\]

$v=2x^3+7x^2+9x+7,\quad w=7x^3+6x$

NWD$(v,w)=1$.

Wyznaczono $g=84.48215863265216$, więc:
\[
w+g=7x^3+6x+84.48215863265216
\]

\[
\mathrm{NWD}(v,w+g)=x+2.1693702699113073
\]

{\small
\begin{align*}
\mathrm{NWW}(v,w+g)= {} & x^5+1.3306x^4+2.4705x^3\\
&+13.2094x^2+17.4421x+19.4716
\end{align*}
}

\section*{3. Porządek produktowy}

\[
\begin{array}{c|c}
\text{Porównanie} & \text{Wynik} \\\hline
(2,7)\le(9,7) & \text{tak}\\
(2,7)\le(0,6) & \text{nie}\\
(9,7)\le(2,7) & \text{nie}\\
(9,7)\le(0,6) & \text{nie}\\
(0,6)\le(2,7) & \text{tak}\\
(0,6)\le(9,7) & \text{tak}
\end{array}
\]

\[
(2,9,0)\not\le(7,7,6),\quad (7,7,6)\not\le(2,9,0)
\]

Minimalne w $A$: $(0,6), (1,5)$.

{\scriptsize
\begin{lstlisting}
[0,0,0,16], [0,0,4,15], [0,0,6,14], [0,0,7,13], [0,0,8,0],
[0,15,0,15], [0,15,4,14], [0,15,6,13], [0,15,7,0],
[0,16,0,14], [0,16,4,13], [0,16,6,0], [0,17,0,13],
[0,17,3,0], [19,0,0,15], [19,0,4,14], [19,0,6,13],
[19,0,7,0], [19,15,0,14], [19,15,4,13], [19,15,5,0],
[19,16,0,13], [19,16,3,0], [19,17,0,0]
\end{lstlisting}
}

\end{document}